\section{Exact conversion: known finite spectra versus intervals}\label{sec:exact}
\begin{proposition}[Two-query conversion]\label{prop:two-point}
Let $0<a<b\leq1$ be known, with $(1-a)(1-b)\leq a+b$. The polynomial
\[
 p_{a,b}(x)=1-\frac{ab+x^2}{a+b}
\]
is bounded in absolute value by one on $[-1,1]$. Two QSVT queries implement it, and
\begin{equation}\label{eq:quadratic-error}
 p_{a,b}(H)-(I-H)=-\frac{(H-aI)(H-bI)}{a+b}.
\end{equation}
Consequently the conversion is exact if $\spec(H)\subseteq\{a,b\}$; for any other spectrum its error is exactly
$\max_{x\in\spec(H)}|(x-a)(x-b)|/(a+b)$.
\end{proposition}
\begin{proof}
The polynomial decreases as $x^2$ increases. Its maximum is $1-ab/(a+b)<1$, and its minimum is $-(1-a)(1-b)/(a+b)\geq-1$. The remaining assertions follow from Lemma~\ref{lem:implementation} and factorization.
\end{proof}
In particular, $b\geq1/2$ suffices for the hypothesis, and $b=1$ gives $p_{a,1}(x)=(1-x^2)/(1+a)$. Thus a large condition number alone does not make unit-normalized conversion to $I-H$ expensive. This result concerns the conversion of an operator; it does not remove the normalization and overlap costs of a linear-system output state.

\begin{proposition}[An interval cannot be converted exactly]\label{prop:no-exact}
No finite-query converter produces an exact encoding of $I-H$ with normalization $0<\nu<2$ for every scalar $H=x$ in any nonempty open subinterval of $(0,1)$.
\end{proposition}
\begin{proof}
By Lemma~\ref{lem:scalar}, the designated output entry $q(t)$ of such a converter is a trigonometric polynomial. By analyticity, equality with $(1-\sin t)/\nu$ on an open interval implies equality for every real $t$. At $t=-\pi/2$ this would give $q(t)=2/\nu>1$, contradicting unitarity.
\end{proof}

\begin{proposition}[Pairwise lower bound and normalization near one]\label{prop:pairwise}
Suppose $0<\delta\leq1/4$ and a converter works at both $H=\delta$ and $H=2\delta$ with error at most $\delta/16$. It uses at least
\[
 \frac{\sqrt3}{2}\left(\sqrt{\frac{93}{32}}-\sqrt{\frac{17}{8}}\right)\delta^{-1/2}
\]
queries. Moreover, the differentiated hybrid argument gives the following formal bound: if $\nu>1-\delta$ and an exact converter with normalization $\nu$ that uses $T$ queries works on a neighborhood of $x=\delta$, then
\begin{equation}\label{eq:normalization-slack}
 T\geq\frac{(1-\delta)\sqrt{1-\delta^2}}
 {\nu\sqrt{\nu^2-(1-\delta)^2}}.
\end{equation}
\end{proposition}
\begin{proof}
Use the canonical oracle $U(\arcsin x)$. Its derivative in $x$ has norm $1/\sqrt{1-x^2}$, so the hybrid argument gives
\[
 \norm{V_{2\delta}-V_\delta}\leq\frac{2T\delta}{\sqrt3}.
\]
Apply each output unitary to its designated input basis vector. If the designated output amplitude is $z$, the component orthogonal to the designated output has norm $\sqrt{1-|z|^2}$. At $x=\delta$ this norm is at most $\sqrt{17\delta/8}$, since $|z|\geq1-17\delta/16$. At $x=2\delta$ it is at least $\sqrt{93\delta/32}$: indeed $|z|\leq1-31\delta/16$, and $1-(1-s)^2\geq3s/2$ for $s=31\delta/16\leq1/2$. The difference between these norms bounds the distance between the output vectors from below, proving the first claim.

For the second claim, the exact designated amplitude is $f(x)=(1-x)/\nu$. The norm of the orthogonal component has derivative $(1-x)/(\nu\sqrt{\nu^2-(1-x)^2})$. Compare neighboring inputs, apply the hybrid argument with oracle derivative $1/\sqrt{1-x^2}$, divide by their separation, and take the limit.
\end{proof}
The pairwise method is a special case of the lower-bound method for eigenvalue transformations in \cite[Theorem 73, extended version]{gilyen2019}. The differentiated inequality \eqref{eq:normalization-slack} is only formal and gives no lower bound: Proposition~\ref{prop:no-exact} rules out its hypothesis for every $\nu<2$, and for $\nu\geq2$ its right side is less than one. It only shows how the local obstruction grows as $\nu$ decreases to $1-\delta$. At normalization two, a direct coherent linear combination implements $(I-H)/2$ with a constant number of queries. The rest of the paper studies how the cost at normalization one depends on the accuracy.
